\subsection{Service primitives and the scope of acceptance}
The motivating inventory contract chooses physical stock and a service tier
before demand is realized. An externally fixed premium is paid by the customer;
shortage indemnity is paid back to that customer. Maintenance and tier changes
consume real resources. The customer values expected fulfillment and transfers,
whereas the provider values receipts net of inventory, shortage, maintenance,
and adjustment costs. A static tier comparator can still choose stock after
observing the public demand regime. Electronic Companion
Sections~\ref{sec:model}--\ref{sec:accepted-continuous} retain the full institutional
model, exact information hierarchy, and all supporting proofs.

Two acceptance institutions must not be conflated. The canonical master
agreement imposes ex ante participation and expected fill commitments. The
continuation institution additionally lets a customer leave at every review.
Once supporting prices identify physical decisions, the latter institution
induces a payment cap for each nonroot subtree and a root payment commitment.
The results below retain every such row; they do not infer continuation
acceptance from ex ante surplus. The benchmark-relative transfer theorem uses
positive scalar payment coefficients. General vector services and signed
coefficients remain covered by the polyhedral balance and certification
results, without applying that scalar coordinate map outside its assumptions.

